\begin{example}
e.g. Let's consider the following task. \vspace{7pt}

\textbf{Task}: multiply two natural numbers $a$ and $b$, both expressed by $n$ digits. \vspace{7pt}

A possible solution could be: reduce the multiplication to a sequence of additions. Therefore, knowing that we have to repeat the addition of number $a$ for $b$
times, our process is:
$$a \times b = a+a+ \dots + a$$
The amount of time required by this process can be summarized as follows:
$$b \times n$$
since we repeat the addition operation for $b$ times. We call this process as \textbf{RepeatedAdditions}. \vspace{7pt}

Another process that can solve the previous task might be: multiply $a$ and $b$ using the basic elementary school algorithm.
Let $a_1, a_2, \dots, a_n$ be the digits of $a$, we compute the final result performing $b \times (a_1 \times 10^0) + \dots + b \times (a_n \times 10^{n-1})$, which
takes a number of steps equal to $n \times n$. We call this process as \textbf{GridMethod}. \vspace{7pt}

Between these two processes there's a huge difference, that can described as follows:
\begin{itemize}
    \renewcommand{\labelitemi}{-}
    \item RepeatedAdditions: the term $b$ can be \textbf{exponential} in $n$, so $b$ might be $10^n$ in the worst case.
    \item GridMethod: its complexity is quadratic, but anyway we are sure that it takes less time or effort than RepeatedAdditions.
\end{itemize}

Even though by both processes we can solve the task, RepeatedAdditions approach takes ages. Therefore, they solve the same task in different ways and not 
always this is acceptable in terms of efficiency.
\end{example}